\documentclass{article}
\usepackage[utf8]{inputenc}
\usepackage[T1]{fontenc}
\usepackage{amsmath, amssymb, amsthm}
\usepackage{geometry}
\usepackage{hyperref}
\usepackage{bookmark}
\usepackage{tikz}

\usepackage{titling}

% \pretitle{\begin{center}\Huge\bfseries}
% \posttitle{\end{center}}
% \preauthor{\begin{center}\large}
% \postauthor{\end{center}}
% \predate{\begin{center}\large}
% \postdate{\end{center}}

% \renewcommand{\thesubsection}{\Roman{subsection}}
% \newcommand{\exo}[2]{\textbf{\Roman{#1} #2} }

% \geometry{left=2.5cm, right=2.5cm, top=2.5cm, bottom=2.5cm}

\newcommand{\subtitle}[1]{
    \posttitle{\par\end{center}\begin{center}\large#1\end{center}\vskip0.5em}
}


\setlength{\parindent}{0pt}
\setlength{\parskip}{0.8em}

\title{Equation différentielle ordinaire}

\subtitle{Methodes de résolution numérique} 

\author{Axel Lavigne}
\date{\today}


\newcommand{\exo}[1]{
    \vspace{5mm}
    \textbf{Exercice #1}
    \hrulefill
}

\newcommand{\smalltext}[1]{
    \scriptsize
    \textit{#1}
    \normalsize
}

\begin{document}

\maketitle

\textbf{Problèmes de Cauchy}

\begin{equation*}
    \frac{dy}{dt} = f(t, y)
\end{equation*}

avec comme conditions initiales:

\begin{equation*}
    y(t_0) = y_0
\end{equation*}


\textbf{Methode d'Euler}

\textit{Explicite}

\begin{align*}
    u_{n+1} &= u_n + h f(t_n, u_n) \\
    t_{n+1} &= t_n + h
\end{align*}

\textit{Implicite}

\begin{align*}
    u_{n+1} &= u_n + h f(t_{n+1}, u_{n+1}) \\
    t_{n+1} &= t_n + h
\end{align*}

\smalltext{avec $h$ le pas de temps, $t_n$ le temps discret et $u_n$ l'approximation de $y(t_n)$}

\smalltext{à noter que la methode implicite est plus stable que la methode explicite, 
mais plus complexe à résoudre pour trouver $u_{n+1}$ en fonction de $u_n$} 

\textbf{Shema}


\smalltext{nous noterons $y(t_n)$ par $y_n$ et $y(t_{n+1})$ par $y_{n+1}$}

Un schéma est une formule \textit{discrétisée} qui transforme une équation différentielle en une relation algébrique

il comporte 3 éléments:

\begin{itemize}
    \item Un temps discret $t_n = t_0 + n h$
    \item Une approximation de la dérivée $y'$ en fonction de $y$ et $t$ : $y'(t_n) \approx \phi(t_n, y_n, h)$
    \item Une relation de récurrence pour calculer $y_{n+1}$ à partir de $y_n$
\end{itemize}

Un shema doit être consistant, stable et convergent pour correctement approximer une solution

\textit{Consistance}: \hrulefill
\smalltext{On note l'erreur de consistance $\epsilon_n$}
\begin{equation*}
    \epsilon_n=y(t_{n+1})-y(t_n)-h_nf(t_n,y(t_n))
\end{equation*}

\smalltext{Un schéma est consistant si $\sum \epsilon_n$ tend vers 0 lorsque $h$ tend vers 0}

% \smalltext{la différence entre la solution exacte et la solution approchée doit tendre vers 0 lorsque $h$ tend vers 0}



\textit{Stabilité}: \hrulefill
\begin{equation*}
    |y(t_n)| \leq C
\end{equation*}
\smalltext{la solution approchée ne doit pas diverger}

\textit{Convergence}: \hrulefill
\begin{equation*}
    |y(t_n) - y_n| \leq C h^p
\end{equation*}
\smalltext{la solution approchée doit converger vers la solution exacte avec un ordre $p$}




\textbf{Methode de Heun}

On substitue $y_{n+1}$ dans la formule Implicite par $y_n + h f(t_n, y_n)$

\begin{align*}
    y_{n+1} &= y_n + \frac{h}{2} (f(t_n, y_n) + f(t_{n+1}, y_n + h f(t_n, y_n))) \\
    t_{n+1} &= t_n + h
\end{align*}



\textbf{Definitions}

\textit{Ordre de convergence}: $p$ tel que $|y(t_n) - y_n| \leq C h^p$

\textit{Stabilité}: la solution approchée ne doit pas diverger

\textit{Consistance}: convergence de la solution approchée vers la solution exacte lorsque $h$ tend vers 0

\textit{Implicite}: $y_{n+1} = y_n + h f(t_{n+1}, y_{n+1})$

\textit{Explicite}: $y_{n+1} = y_n + h f(t_n, y_n)$

\textit{Schéma}: formule discrétisée pour résoudre une équation différentielle

\textit{discrétisée}: rendre un problème continu en un problème \textit{discret}

\textit{discret}: séparé en éléments distincts



\end{document}


